\documentclass[12pt,a4paper]{article}
\usepackage[utf8]{inputenc}
\usepackage[spanish]{babel}
\usepackage[margin=2.5cm]{geometry}
\usepackage{graphicx}
\graphicspath{{imagenes/}}
\usepackage{amsmath}
\usepackage{amsfonts}
\usepackage{amssymb}
\usepackage{booktabs}
\usepackage{hyperref}
\usepackage{float}

\begin{document}

% --- PORTADA ---
\begin{titlepage}
    \centering
    {\Large \textbf{Universidad Distrital Francisco José de Caldas} \\[0.3cm]}
    {\large Proyecto Curricular de Ingeniería en Sistemas \\[0.2cm]}
    {\large Curso de Arquitectura de Computadores \\[1.5cm]}
    
    {\huge \textbf{Práctica 3: Construcción de una ALU de 8 Bits con Circuitos Combinacionales} \\[2cm]}
    
    \textbf{Docente:} \\
    Ing. Jonathan Roberto Torres Castillo, PhD. \\[1cm]
    
    \textbf{Estudiantes:} \\
    Juan Camilo Carvajal Camargo \\
    Jenny Vanesa Leon Pinto \\[2cm]
    
    \vfill
    {\large 03 de octubre de 2026}
\end{titlepage}

\newpage

% --- ÍNDICE ---
\tableofcontents
\newpage

% --- 1. INTRODUCCIÓN ---
\section{Introducción}
La Unidad Aritmético-Lógica (ALU) constituye un componente fundamental en la arquitectura de los sistemas digitales[cite: 1], al encargarse de la ejecución de operaciones aritméticas y lógicas sobre datos binarios. Su diseño se basa en la integración de circuitos combinacionales que implementan funciones booleanas específicas, permitiendo la realización de operaciones como suma, resta y operaciones lógicas básicas (AND, OR, NOT, entre otras)[cite: 1].

Desde una perspectiva de diseño digital, la construcción de una ALU implica la descomposición de operaciones complejas en bloques funcionales elementales, los cuales pueden ser implementados mediante compuertas lógicas[cite: 1, 2]. Este enfoque permite comprender la relación entre las expresiones algebraicas booleanas, su representación mediante tablas de verdad y su materialización en circuitos físicos[cite: 2].

En esta práctica se aborda el diseño e implementación de una ALU de 8 bits utilizando exclusivamente compuertas lógicas y circuitos combinacionales, sin recurrir a bloques predefinidos[cite: 2]. Esto permite analizar de manera estructurada la construcción de sumadores, unidades lógicas y mecanismos de selección de operaciones, integrándolos en un sistema funcional que emula el comportamiento básico de una ALU[cite: 2].

% --- 2. OBJETIVOS ---
\section{Objetivos}
\begin{itemize}
    \item Comprender la estructura funcional de una ALU de 8 bits[cite: 2].
    \item Diseñar y construir circuitos combinacionales para operaciones aritmético-lógicas[cite: 2].
    \item Implementar un multiplexor de 8 a 1 de 8 bits sin componentes predefinidos[cite: 2].
    \item Integrar todos los módulos en un circuito completo de ALU[cite: 2].
    \item Evaluar el funcionamiento de la ALU mediante casos de prueba[cite: 2].
\end{itemize}

% --- 3. MARCO TEÓRICO ---
\section{Marco Teórico}
Los sistemas digitales se fundamentan en la manipulación de información binaria mediante circuitos combinacionales, los cuales implementan funciones booleanas a partir de compuertas lógicas[cite: 2]. 

\subsection{Unidad Aritmético-Lógica (ALU)}
La ALU es un bloque funcional encargado de ejecutar operaciones aritméticas y lógicas sobre datos binarios[cite: 2]. Su operación se basa en la recepción de operandos y una señal de control, la cual determina la operación que se debe realizar[cite: 2]. Es un circuito de tipo combinacional, lo que implica que sus salidas dependen únicamente de los valores actuales de sus entradas[cite: 2].

\subsection{Operadores Lógicos}
Los operadores lógicos permiten realizar transformaciones directas sobre variables binarias mediante reglas del álgebra de Boole, aplicándose bit a bit[cite: 3]:
\begin{itemize}
    \item \textbf{AND:} $Y_i = A_i \cdot B_i$[cite: 3]
    \item \textbf{OR:} $Y_i = A_i + B_i$[cite: 3]
    \item \textbf{XOR:} $Y_i = A_i \overline{B_i} + \overline{A_i} B_i$[cite: 3]
    \item \textbf{NOT:} $Y_i = \overline{A_i}$[cite: 3]
    \item \textbf{XNOR:} $Y_i = A_i B_i + \overline{A_i} \overline{B_i}$[cite: 3]
\end{itemize}

\subsection{Operadores Aritméticos}
\begin{itemize}
    \item \textbf{Sumador completo (Full Adder):} Circuito combinacional que permite sumar tres bits de entrada (dos operandos y un acarreo de entrada $C_{in}$), produciendo un bit de suma ($Sum$) y un acarreo de salida ($C_{out}$)[cite: 3].
    \begin{equation}
        Sum = A \oplus B \oplus C_{in}
    \end{equation}
    \begin{equation}
        C_{out} = AB + C_{in}(A \oplus B)
    \end{equation}
    \item \textbf{Resta en binario:} Se basa en el uso del complemento a dos, transformando la resta en una suma[cite: 3]:
    \begin{equation}
        A - B = A + (\overline{B} + 1)
    \end{equation}
    \item \textbf{Multiplicación binaria:} Se fundamenta en la generación de productos parciales mediante compuertas AND y su posterior suma con desplazamientos[cite: 4].
\end{itemize}

\subsection{Multiplexor de 8 a 1 de 8 bits}
Un multiplexor selecciona una entre 8 entradas utilizando 3 líneas de control $(S_2, S_1, S_0)$[cite: 6]. Para manejar una salida de 8 bits, se conectan 8 multiplexores en paralelo, compartiendo las mismas líneas de selección[cite: 6].

% --- 4. DESCRIPCIÓN DEL DESARROLLO ---
\section{Descripción del Desarrollo}
En esta sección se detallan los pasos seguidos para el diseño e implementación de cada módulo en \textbf{Logisim-Evolution}[cite: 6]:
\begin{enumerate}
    \item \textbf{Componentes Básicos (Operaciones Lógicas):} Implementación mediante compuertas AND, OR, XOR, NOT y XNOR dispuestas en paralelo para buses de 8 bits[cite: 7].
    \item \textbf{Sumador y Restador de 8 bits:} Encadenamiento en configuración \textit{ripple-carry} de 8 sumadores completos de 1 bit[cite: 7]. Para la resta, se configuró la inversión de $B$ con $C_{in} = 1$[cite: 8].
    \item \textbf{Multiplicador de 8 bits:} Diseñado mediante productos parciales y desplazamientos estructurales[cite: 8].
    \item \textbf{Multiplexor e Integración:} Se acoplaron las salidas de las operaciones al MUX de 8 a 1 controlado por el Opcode $(S_2 S_1 S_0)$[cite: 8].
\end{enumerate}

% --- 5. RESULTADOS ---
\section{Resultados}
A continuación se relacionan las operaciones de la ALU según el Opcode[cite: 9]:

\begin{table}[H]
    \centering
    \begin{tabular}{lc}
        \toprule
        \textbf{Operación} & \textbf{Opcode $(S_2 S_1 S_0)$} \\
        \midrule
        AND & 000 \\
        OR & 001 \\
        XOR & 010 \\
        NOT & 011 \\
        XNOR & 100 \\
        Resta & 101 \\
        Suma & 110 \\
        Multiplicación & 111 \\
        \bottomrule
    \end{tabular}
    \caption{Codificación del Opcode para la ALU[cite: 9]}
    \label{tab:opcode}
\end{table}

\noindent Se verificó exitosamente el funcionamiento de la simulación. Por ejemplo, al realizar una multiplicación con las entradas $A = 00000011_2$ ($3_{10}$) y $B = 00000100_2$ ($4_{10}$) bajo el Opcode $111$, se obtuvo la salida correcta $F = 00001100_2$ ($12_{10}$)[cite: 9].

% --- 6. ANÁLISIS Y DISCUSIÓN ---
\section{Análisis y Discusión}
El diseño modular implementado permitió aislar posibles fallas en cada bloque antes de la integración final con el multiplexor[cite: 7, 8]. Una de las principales dificultades técnicas radicó en la correcta propagación del acarreo en el sumador/restador de 8 bits y en la gestión de los desplazamientos para la multiplicación[cite: 7, 8]. No obstante, el uso de circuitos combinacionales puros garantizó respuestas síncronas estables sin retardos secuenciales indeseados.

% --- 7. CONCLUSIONES ---
\section{Conclusiones}
\begin{itemize}
    \item La implementación jerárquica demostró que es posible construir sistemas complejos de procesamiento aritmético-lógico partiendo de compuertas básicas elementales[cite: 2, 7].
    \item El uso del complemento a dos optimizó los recursos de hardware al reutilizar la estructura del sumador de 8 bits para ejecutar operaciones de resta[cite: 3, 8].
    \item La arquitectura basada en un multiplexor de 8 a 1 permitió consolidar múltiples rutas de datos independientes en una única salida controlada mediante líneas de selección estables[cite: 6, 8].
\end{itemize}

% --- 8. ENLACE A CARPETA ---
\section{Enlace a la Carpeta Compartida}
\url{https://github.com/} \textit{(Reemplazar con el enlace real al archivo .circ en Logisim)}[cite: 10, 11]

% --- 9. REFERENCIAS ---
\section{Referencias}
\begin{thebibliography}{9}
    \bibitem{mano2017} Mano, M. M., \& Ciletti, M. D. (2017). \textit{Digital Design} (6th ed.). Pearson[cite: 13].
    \bibitem{floyd2019} Floyd, T. L. (2019). \textit{Digital Fundamentals} (11th ed.). Pearson[cite: 13].
    \bibitem{roth2015} Roth, C. H., \& Kinney, L. L. (2015). \textit{Fundamentals of Logic Design} (7th ed.). Cengage Learning[cite: 13].
    \bibitem{logisim} Logisim Evolution. (n.d.). \textit{Open-source digital logic simulator}. Recuperado de \url{https://github.com/logisim-evolution/logisim-evolution}[cite: 13].
\end{thebibliography}
LaboratorioArquiIII
\end{document}